%======================================================================
% PRESENTACIÓN
%======================================================================
\chapter*{Presentación}
\addcontentsline{toc}{chapter}{Presentación}

Sandra Ecosystem representa una propuesta integral para la transformación digital de instituciones gubernamentales y empresariales en la República Bolivariana de Venezuela. Este documento técnico tiene como propósito principal presentar de manera exhaustiva la arquitectura, componentes, funcionalidades y estándares que rigen el ecosistema, proporcionando una guía completa para comprender su diseño y potencial de implementación.

El documento está estructurado siguiendo los cánones de la documentación técnica académica, comenzando con un marco teórico que contextualiza el proyecto dentro del panorama tecnológico actual, seguido de un marco legal que fundamenta la propuesta desde la perspectiva normativa venezolana. Posteriormente, se presenta la arquitectura general del ecosistema y el detalle de cada uno de sus componentes principales.

La relevancia de este trabajo radica en la necesidad creciente de nuestras instituciones de contar con soluciones tecnológicas propias, capaces de adaptarse a las particularidades del contexto nacional mientras mantienen los más altos estándares de calidad, seguridad y rendimiento. Sandra Ecosystem responde a esta necesidad ofreciendo una plataforma robusta, escalable y completamente adaptable.

Este documento está dirigido a profesionales de tecnología, directivos de instituciones, arquitectos de sistemas y todos aquellos interesados en comprender las capacidades y el potencial del ecosistema. Se espera que el lector adquiera una comprensión integral de la plataforma que le permita evaluar su aplicabilidad en diferentes contextos organizacionales.

\newpage
